% GENERATED FILE. Edit canonical lessons, then run npm run generate:books.
% canonical-source-hash: fnv1a64:2a11ae5176c101e7
% canonical-lessons: GE-C396-tastatur, GE-C396-drucker, GE-C396-ausdrucken, GE-C396-installieren, GE-C396-programm

\chapter{Am Computer}
\label{ch:ge-am-computer}
\hlchaptermodality{396}

\begin{chapteropening}
\textbf{By the end of this chapter:} \emph{I can talk about the keyboard and the printer, printing something out, and installing a program.}

Complete the last lesson of chapter 396: I can talk about the keyboard and the printer, printing something out, and installing a program.
\end{chapteropening}

\section[die Tastatur]{die Tastatur — the keyboard}

\label{lesson:GE-C396-tastatur}



\begin{quote}
Before the new one: say the German for data protection, then the German for to encrypt.
\end{quote}

\subsection*{You'll want to know: die Tastatur}
\textbf{die Tastatur} — \textquotedblleft{}the keyboard\textquotedblright{}.

\begin{quote}
Auf der deutschen Tastatur sind Y und Z vertauscht.
\end{quote}

\begin{quote}
\emph{On the German keyboard Y and Z have swapped places.}
\end{quote}

Plural: \emph{die Tastaturen}. \emph{Die Taste} — the key.

\begin{cousinweb}
\emph{Tastatur} is built on \emph{Taste} \textquotedblleft{}key\textquotedblright{}, borrowed from Italian \emph{tasto} \textquotedblleft{}key\textquotedblright{}, from \emph{tastare} \textquotedblleft{}to touch, to feel\textquotedblright{}, which goes back to Latin \emph{taxitāre} (EtymDB), with the Latin-style ending \emph{-atur}.
\end{cousinweb}

\begin{grammarlens}[title={What you've built}]
One more word for telling, explaining and describing at B1.
\end{grammarlens}

\subsection*{Your turn}
\begin{itemize}
\raggedright
  \item \emph{Say it:} \emph{die Tastatur}
  \item \emph{Say it:} \emph{die Tastatur}, once more
  \item \emph{Say it:} a sentence of your own with \emph{die Tastatur}
  \item \emph{Recall:} say the German for data protection, then the German for to encrypt, then say \emph{die Tastatur} again
\end{itemize}

\begin{tcolorbox}[breakable,colback=teal!4,colframe=teal!35!black,title={Before you move on}]
What does die Tastatur mean? (\textquotedblleft{}The keyboard\textquotedblright{}.) Say it once more.
\end{tcolorbox}
\section[der Drucker]{der Drucker — the printer}

\label{lesson:GE-C396-drucker}



\begin{quote}
Before the new one: say the German for to encrypt, then the German for the keyboard.
\end{quote}

\subsection*{You'll want to know: der Drucker}
\textbf{der Drucker} — \textquotedblleft{}the printer\textquotedblright{}.

\begin{quote}
Der Drucker hat kein Papier mehr.
\end{quote}

\begin{quote}
\emph{The printer has run out of paper.}
\end{quote}

Plural: \emph{die Drucker}. Also a printer by trade.

\begin{cousinweb}
\emph{Drucker} is formed with \emph{-er} on \emph{drucken} \textquotedblleft{}to print\textquotedblright{}, of the family of \emph{drücken} \textquotedblleft{}to press\textquotedblright{}, from Proto-Germanic \emph{þrukjaną} \textquotedblleft{}to press\textquotedblright{} (EtymDB).
\end{cousinweb}

\begin{grammarlens}[title={What you've built}]
One more word for telling, explaining and describing at B1.
\end{grammarlens}

\subsection*{Your turn}
\begin{itemize}
\raggedright
  \item \emph{Say it:} \emph{der Drucker}
  \item \emph{Say it:} \emph{der Drucker}, once more
  \item \emph{Say it:} a sentence of your own with \emph{der Drucker}
  \item \emph{Recall:} say the German for to encrypt, then the German for the keyboard, then say \emph{der Drucker} again
\end{itemize}

\begin{tcolorbox}[breakable,colback=teal!4,colframe=teal!35!black,title={Before you move on}]
What does der Drucker mean? (\textquotedblleft{}The printer\textquotedblright{}.) Say it once more.
\end{tcolorbox}
\section[ausdrucken]{ausdrucken — to print out}

\label{lesson:GE-C396-ausdrucken}



\begin{quote}
Before the new one: say the German for the keyboard, then the German for the printer.
\end{quote}

\subsection*{You'll want to know: ausdrucken}
\textbf{ausdrucken} — \textquotedblleft{}to print out\textquotedblright{}.

\begin{quote}
Ich drucke die Fahrkarte zu Hause aus.
\end{quote}

\begin{quote}
\emph{I print out the ticket at home.}
\end{quote}

Separable: \emph{ich drucke aus}, \emph{ich habe ausgedruckt}. Careful: \emph{ausdrücken}, with an umlaut, means to express.

\begin{cousinweb}
\emph{ausdrucken} is built from \emph{aus} \textquotedblleft{}out\textquotedblright{} and \emph{drucken} \textquotedblleft{}to print\textquotedblright{}, from Proto-Germanic \emph{þrukjaną} \textquotedblleft{}to press\textquotedblright{} (EtymDB): you press the text out onto paper.
\end{cousinweb}

\begin{grammarlens}[title={What you've built}]
One more word for telling, explaining and describing at B1.
\end{grammarlens}

\subsection*{Your turn}
\begin{itemize}
\raggedright
  \item \emph{Say it:} \emph{ausdrucken}
  \item \emph{Say it:} \emph{ausdrucken}, once more
  \item \emph{Say it:} a sentence of your own with \emph{ausdrucken}
  \item \emph{Recall:} say the German for the keyboard, then the German for the printer, then say \emph{ausdrucken} again
\end{itemize}

\begin{tcolorbox}[breakable,colback=teal!4,colframe=teal!35!black,title={Before you move on}]
What does ausdrucken mean? (\textquotedblleft{}To print out\textquotedblright{}.) Say it once more.
\end{tcolorbox}
\section[installieren]{installieren — to install}

\label{lesson:GE-C396-installieren}



\begin{quote}
Before the new one: say the German for the printer, then the German for to print out.
\end{quote}

\subsection*{You'll want to know: installieren}
\textbf{installieren} — \textquotedblleft{}to install\textquotedblright{}.

\begin{quote}
Mein Sohn hat die App auf meinem Handy installiert.
\end{quote}

\begin{quote}
\emph{My son installed the app on my mobile.}
\end{quote}

Perfekt: \emph{hat installiert}. Also of wiring: \emph{eine Lampe installieren}.

\begin{cousinweb}
\emph{installieren} is a learned borrowing from Medieval Latin \emph{installāre} (EtymDB), built on \emph{stallum} \textquotedblleft{}a stall, a seat\textquotedblright{}: usually said to have meant first to place someone in a church seat or an office.
\end{cousinweb}

\begin{grammarlens}[title={What you've built}]
One more word for telling, explaining and describing at B1.
\end{grammarlens}

\subsection*{Your turn}
\begin{itemize}
\raggedright
  \item \emph{Say it:} \emph{installieren}
  \item \emph{Say it:} \emph{installieren}, once more
  \item \emph{Say it:} a sentence of your own with \emph{installieren}
  \item \emph{Recall:} say the German for the printer, then the German for to print out, then say \emph{installieren} again
\end{itemize}

\begin{tcolorbox}[breakable,colback=teal!4,colframe=teal!35!black,title={Before you move on}]
What does installieren mean? (\textquotedblleft{}To install\textquotedblright{}.) Say it once more.
\end{tcolorbox}
\section[das Programm]{das Programm — the program; the programme}

\label{lesson:GE-C396-programm}



\begin{quote}
Before the new one: say the German for to print out, then the German for to install.
\end{quote}

\subsection*{You'll want to know: das Programm}
\textbf{das Programm} — \textquotedblleft{}the program; the programme\textquotedblright{}.

\begin{quote}
Mit diesem Programm kann man Fotos bearbeiten.
\end{quote}

\begin{quote}
\emph{You can edit photos with this program.}
\end{quote}

Plural: \emph{die Programme}. Also a TV channel or a schedule: \emph{das Fernsehprogramm}.

\begin{cousinweb}
\emph{Programm} is a learned borrowing from Late Latin \emph{programma} \textquotedblleft{}a proclamation, an edict\textquotedblright{} (Lewis \& Short), from Greek \emph{prógramma} \textquotedblleft{}a public notice\textquotedblright{}, built from \emph{pro} \textquotedblleft{}before\textquotedblright{} and \emph{gráphein} (EtymDB).
\end{cousinweb}

\begin{grammarlens}[title={What you've built}]
One more word for telling, explaining and describing at B1.
\end{grammarlens}

\subsection*{Your turn}
\begin{itemize}
\raggedright
  \item \emph{Say it:} \emph{das Programm}
  \item \emph{Say it:} \emph{das Programm}, once more
  \item \emph{Say it:} a sentence of your own with \emph{das Programm}
  \item \emph{Recall:} say the German for to print out, then the German for to install, then say \emph{das Programm} again
\end{itemize}

\begin{tcolorbox}[breakable,colback=teal!4,colframe=teal!35!black,title={Before you move on}]
What does das Programm mean? (\textquotedblleft{}The program; the programme\textquotedblright{}.) Say it once more.
\end{tcolorbox}
